\subsection{One connector from the terminal \((1,2)\)}
\label{proof:Q5-C032}

Use the original finite measure and a real skew kernel \(K\), \(|K|\le1\).
Write \(A=J+K\), \(T=J-K\), \(g=K\one\), \(p=1+g\),
\(q=1-g\), \(f_2=A^2\one\), \(d=\E g^2\),
\(\delta_j=\|K^j\one\|^2\), and \(\theta=13/32\).
Lemma~\ref{q3:lem:spectral} gives \(0\le d\le1/3\),
\(\|K\|^2\le\theta\), and
\(\delta_{j+1}\le\theta\delta_j\).
Direct the terminal unit arm inward, its two-edge arm outward, the
connector toward the receiving root, and that root's length-two arm
outward. Thus \(F=qf_2\) and \(X=TF\). We prove
\begin{equation}\tag{L32.T}\label{eq:L32:claim}
 \E[f_2X^4]\le1.
\end{equation}
The inequality is strict for \(d>0\). Both directions and the reference
weight are fixed throughout the proof.

\paragraph{The common covariance criterion.}
Put \(r=Kg\), \(s=Kr\), \(U=g(g+r)\),
\(m=1-2d\), \(k=1-d\), and \(h=4m^3\).
Since \(f_2=1+g+r\in[0,2]\), \(|g+r|\le1\). Skewness gives
\[
 \E U=d,\quad \E U^2\le d,\quad
 \operatorname{Var}U\le dk,\quad
 \E|U|\le\sqrt{d(d+\delta_2)}.
\]
The last bound is Cauchy--Schwarz; hence
\(\|KU\|^2\le d^2(1+v)\), where \(v=\delta_2/d\) when \(d>0\).
The actual identities, with \(Z=X-m\), are
\begin{equation}\tag{L30.2}\label{eq:L30:ledger}
\begin{aligned}
 F&=1+r-U,\qquad X=q-2d-s+KU,\\
 Z&=-g-s+KU,\qquad \E f_2=k,\\
 \E[f_2Z]&=-d+\delta_2-\E[(g+r+s)U],\\
 \E Z^2&=d-2\delta_2+\delta_3+\|KU\|^2\\
 &\quad+2\E[(r+K^2r)U].
\end{aligned}
\end{equation}
Here \(m=\E F\); it need not equal \(\E X\).
Let \(c\ge0\) satisfy the actual pointwise bound
\(f_2(X^2+2mX+3m^2)\le c\).
Expanding the fourth power at \(m\), and using
\(\delta_3\le\theta\delta_2\), gives
\[
 \E[f_2X^4]-1\le-d\mathcal L
                  +\E[(V-\E V)(U-d)],
\]
where
\begin{align*}
 V&=-h(g+s)+(2c-h)r+2cK^2r,\\
 \mathcal B&=3m^4+4m^3+2m^2+2m+2,\\
 \mathcal L&=\mathcal B-kc+[(2-\theta-3d)c-h]v,\\
 \E V&=d(h-2c+2cv).
\end{align*}
The uncentered constant before subtracting one is
\(km^4-hd+h\delta_2+c[d-2\delta_2+\theta\delta_2+d^2(1+v)]
+d^2(h-2c+2cv)\); it equals \(-d\mathcal L+1\).

Let \(\sigma\) be the spectral measure of \(-K^2\) at \(g\).
It has mass \(d\), support \([0,\theta]\) and moments
\(\int x^j\,d\sigma=\delta_{j+1}\). Thus
\[
 \|V\|^2=\int\{h^2(1-x)^2+x(2c-h-2cx)^2\}\,d\sigma.
\]
If \(b=4c^2(2-\theta)-4ch-h^2\ge0\), use
\(\delta_4\le\theta\delta_3\) and the integrated restriction
\(\delta_3\ge\delta_2^2/d\), then subtract \((\E V)^2\):
\begin{equation}\tag{L30.3.3}\label{eq:L30:Q}
 \operatorname{Var}V\le d\mathcal Q,
 \qquad
 \mathcal Q=h^2+(4c^2-4ch-h^2)v-bv^2-d(h-2c+2cv)^2.
\end{equation}
In particular \(\mathcal Q\ge0\) for every actual kernel to which this
calculation applies. Original-measure Cauchy--Schwarz now gives
\begin{equation}\tag{L30.3.5}\label{eq:L30:sharedbound}
 \E[f_2X^4]\le1+d\{-\mathcal L+\sqrt{k\mathcal Q}\}.
\end{equation}
Consequently it suffices to prove
\(b\ge0\), \(\mathcal L>0\), and
\(\mathcal D:=\mathcal L^2-k\mathcal Q>0\).

\paragraph{One certificate on the whole range \(0<d\le1/5\).}
The proved pointwise coefficient \eqref{eq:L29:claim} supplies
\begin{equation}\label{eq:reference-synthesis:coefficient}
 c=255/32-13d+2(v-3/10)_+.
\end{equation}
The actual moment range is \(d\le v\le\theta\), since
\(d^2\le\delta_2\le\theta d\).
Here \(859/160\le c\le1309/160\), \(0<h\le4\); hence
\[
 b\ge4(51/32)(859/160)^2-16(859/160)-16
       =16762811/204800>80.
\]
Here the lower bound first uses \(h\le4\); the resulting polynomial
\(4(51/32)c^2-16c-16\) increases for \(c\ge859/160\).
Also \([(2-\theta-3d)c-h]\ge34181/25600>4/3\).
Writing \(m\in[3/5,1]\), one obtains
\[
 \mathcal L\ge P(m):=
 3m^4+4m^3-\tfrac54m^2-\tfrac{669}{320}m+\tfrac{371}{320}.
\]
Indeed \(c\le1309/160-13d\). Since
\(P''=36m^2+24m-5/2>0\) and
\(P'(3/5)=26571/8000>0\),
\(\mathcal L\ge P(3/5)=3539/5000>0\).

It remains to check a single polynomial on each side of \(v=3/10\).
An expanded form useful for this check is
\[
\begin{aligned}
 \mathcal Q={}&h^2(k-v+v^2)
       +ch[4d-4(1+d)v+4v^2]\\
       &+c^2[-4d+4(1+2d)v-4(2-\theta+d)v^2].
\end{aligned}
\]
For the lower branch put \(v=d+(3/10-d)t\); for the upper branch
put \(v=3/10+(\theta-3/10)t\), with \(0\le t\le1\).
The latter entire rectangle is allowed because \(d\le1/5<3/10\).
These domains retain the actual restriction \(v\ge d\), but need
not consist entirely of realizable kernel moments.

For \(i=0,1,2,3\), \(j=0,1\), substitute
\(d=(i+x)/20\), \(t=(j+y)/2\), \((x,y)\in[0,1]^2\).
The resulting polynomial \(\mathcal D_{ij}\) has degree at most
\((n,m)=(9,2)\) on the lower branch and \((8,4)\) on the upper branch.
Apply Lemma~\ref{lem:bernstein-certificate} to its power coefficients.
The Bernstein coefficients \(\beta_{rs}\) are the explicit rationals
given by \eqref{eq:bernstein-coefficients}.
Every such coefficient is at least the corresponding entry below.
The two columns in each branch correspond to \(j=0,1\).
\begin{table}[htbp]\centering\small
\caption{Lower bounds for every Bernstein coefficient in the sixteen cells.}
\label{tab:reference-synthesis:bernstein}
\begin{tabular}{c|cc|cc}
\(i\)&\multicolumn{2}{c|}{lower branch}&\multicolumn{2}{c}{upper branch}\\
\hline
0&3&5&13&18\\
1&\(3/2\)&\(3/2\)&4&\(13/2\)\\
2&\(3/5\)&\(3/40\)&\(3/4\)&\(5/3\)\\
3&\(1/2\)&\(1/200\)&\(1/25\)&\(7/40\)
\end{tabular}
\end{table}
This finite table follows by expanding the displayed \(\mathcal L\)
and \(\mathcal Q\) and applying
\eqref{eq:bernstein-coefficients}. 
All denominators are positive integers; none depends on \(d,v,x,y\).
The Bernstein basis is nonnegative and sums to one on the unit square.
Thus \(\mathcal D\ge1/200>0\) on both closed branches, including
their common boundary. This proves \eqref{eq:L32:claim} for
\(0<d\le1/5\) by \eqref{eq:L30:sharedbound}.
Indeed \(\mathcal L\le13+(51/32)(1309/160)\theta<19\), and
\(\mathcal L-\sqrt{k\mathcal Q}
=\mathcal D/(\mathcal L+\sqrt{k\mathcal Q})\ge1/7600\).
Thus this range has deficit at least \(d/7600\).

\paragraph{The same direction on \(1/5\le d\le1/3\).}
The direct actual first moment is
\[
 \E[f_2X]=1-4d+\delta_2+2d^2-\E[(g+r+s)U].
\]
Since \(g+s\) is orthogonal to \(r\),
\(\|g+r+s\|^2=d-\delta_2+\delta_3
 \le d-(1-\theta)\delta_2\), and \(\|U\|^2\le d\).
Writing \(\kappa=19/32\), this gives
\[
 \E[f_2X]\le1-4d+2d^2+d[v+\sqrt{1-\kappa v}]
 \le1-\tfrac{87}{32}d+2d^2
 \le h_*(d):=1-\tfrac{81}{32}d+2d^2.
\]
Here \(v+\sqrt{1-\kappa v}\) increases on \([0,\theta]\),
and its right endpoint is less than \(13/32+7/8=41/32\).
The original positive row gives \(0\le X\le m_*(d)=2(1-2d)\).
The functions \(m_*\) and \(h_*\) are nonnegative and decreasing
on this interval; therefore
\[
 \E[f_2X^4]\le m_*^3h_*
 \le(6/5)^3(459/800)=12393/12500<1.
\]
The two parameter ranges use the same original orientation and meet
at \(d=1/5\). At \(d=0\), \(K\one=0\), so every actual tree message
is one by leaf integration even if \(K\ne0\).
For actual equal-pair factors \(0\le E_i\le1\), positivity gives
\(T((\prod_iE_i)qf_2)\le X\), and a receiving equal-pair factor
is also at most one. Full reversal sends \(K\) to \(-K\) and gives
\(\E[\bar f_2\{A(p\bar f_2)\}^4]\le1\) with the same strictness.
